\chapter{Quadratic Forms and the Local--Global Principle}
\begin{center}\itshape
Does $2x^2+7y^2=3z^2$ have a solution in rational numbers, other than $0$?\\
It has real solutions, and $3$-adic ones, and $5$-adic ones. It fails in $\Q_2$ and in $\Q_7$, so the answer is no.\\
And it fails at exactly two places, because failures always come in pairs.
\end{center}

On 10 July 1796 the nineteen-year-old Gauss wrote in his diary, in capital Greek letters, $\text{E}\Upsilon\text{PHKA}$:
$\text{num}=\Delta+\Delta+\Delta$. Every positive integer is a sum of three triangular numbers. Equivalently, every
$8n+3$ is a sum of three odd squares, and behind it lies the theorem that an integer is a sum of three squares unless it
has the form $4^a(8b+7)$. Legendre had found the statement a few years earlier and published a ``proof'' in 1798, which
assumed that every arithmetic progression $an+b$ with $\gcd(a,b)=1$ contains a prime. That was not known; Dirichlet proved
it only in 1837. Gauss's own proof, in the \emph{Disquisitiones} of 1801, used a heavy theory of ternary quadratic forms that
few readers ever followed to the end.

A century later the story turned. Hermann Minkowski, in 1890, classified quadratic forms over $\Q$ by congruence conditions,
in a language so difficult that it was hardly noticed. In 1921 Helmut Hasse, the student who had found Hensel's book in a
bookshop (Chapter~6), proved in his thesis the principle that bears his name: \emph{a quadratic form with rational
coefficients has a non-trivial rational zero if and only if it has one in $\R$ and in every $\Q_p$.} The local conditions
are finitely many and can be checked mechanically, so every question of the form ``does this quadratic equation have a
rational solution?'' becomes a finite computation. The three-squares theorem then takes a few lines, once a clever
geometric trick, found by L\'eon Aubry in 1912 and rediscovered by Davenport and Cassels, turns rational solutions into
integer ones.

The tool that makes the local conditions computable is the \emph{Hilbert symbol} $(a,b)_v=\pm1$, which says whether
$z^2=ax^2+by^2$ has a non-zero solution in $\Q_v$. Its formulas are explicit, and its product over all places is $1$:
that is Hilbert's reformulation of quadratic reciprocity, and it is why failures come in pairs. We follow Serre's
\emph{Course in Arithmetic}, which remains the most elegant road. The proof of the local--global principle goes by the
number of variables: two is easy, three is a descent in the style of Fermat, four needs Dirichlet's theorem on primes in
progressions (so Legendre's missing piece comes back), and five or more reduce to four.

Is all this special to degree $2$? Yes, dramatically so, and we will see why at the end.

\begin{reminder}[from Tomes 1 and 5, and Chapter 6]
\emph{Quadratic reciprocity} (Tome~1, Theorem~6.5): for distinct odd primes, $\bigl(\frac pq\bigr)\bigl(\frac qp\bigr)=
(-1)^{\frac{p-1}2\frac{q-1}2}$, with $\bigl(\frac{-1}p\bigr)=(-1)^{\frac{p-1}2}$ and $\bigl(\frac2p\bigr)=(-1)^{\frac{p^2-1}8}$.
\emph{Squares in $\Q_p$} (Tome~1, Corollary~12.8): $p^vu$ is a square iff $v$ is even and $u$ is a square modulo $p$ ($p$ odd)
or $u\equiv1\pmod8$ ($p=2$). Hence $\Q_p^\times/(\Q_p^\times)^2$ has representatives $1,u,p,up$ ($p$ odd, $u$ a non-residue) and
$\pm1,\pm2,\pm5,\pm10$ ($p=2$), and the squares form an open subgroup. \emph{Dirichlet's theorem} (Tome~5, Theorem~4.1): every
progression $mk+a$ with $\gcd(a,m)=1$ contains infinitely many primes. We write $\Q_\infty=\R$, and $v$ runs over the
\emph{places} $2,3,5,\dots,\infty$.
\end{reminder}

\section{Quadratic forms over a field}

Let $K$ be a field of characteristic $\ne2$. A \emph{quadratic form} in $n$ variables is $f(x)=\sum_{i,j}a_{ij}x_ix_j$ with
$a_{ij}=a_{ji}\in K$; its \emph{bilinear form} is $B(x,y)=\sum a_{ij}x_iy_j$, so that $f(x)=B(x,x)$ and $B(x,y)=\frac12\bigl(f(x+y)
-f(x)-f(y)\bigr)$. The form is \emph{non-degenerate} if $\det(a_{ij})\ne0$, which we always assume. Two forms are
\emph{equivalent}, $f\cong g$, if $g(x)=f(Ax)$ for some $A\in GL_n(K)$; then $\det(g)=\det(A)^2\det(f)$, so the
\emph{discriminant} $d(f)=\det(a_{ij})\in K^\times/K^{\times2}$ is an invariant. We write $\langle a_1,\dots,a_n\rangle$ for
$a_1x_1^2+\dots+a_nx_n^2$, and $f\oplus g$ for the sum of two forms in separate variables.

$f$ \emph{represents} $a\in K$ if $f(x)=a$ for some $x\ne0$, and $f$ is \emph{isotropic} if it represents $0$.

\begin{proposition}\label{prop:diag}
\begin{enumerate}[label=(\alph*)]
\item Every form is equivalent to a diagonal form $\langle a_1,\dots,a_n\rangle$. If $f$ represents $a\ne0$, we may take
$a_1=a$.
\item If $f$ is isotropic, then $f\cong\langle1,-1\rangle\oplus g$ for some $g$, and $f$ represents every element of $K$.
\item Let $a\ne0$. Then $f$ represents $a$ if and only if $f\oplus\langle-a\rangle$ is isotropic.
\end{enumerate}
\end{proposition}
\begin{proof}
Think of $f$ as a function on $V=K^n$. (a) Let $e_1$ be a vector with $f(e_1)=a\ne0$ (one exists unless $f=0$, since $B$ is
recovered from $f$). Every $x\in V$ is $\frac{B(x,e_1)}a e_1+\bigl(x-\frac{B(x,e_1)}ae_1\bigr)$ with the second part in
$e_1^\perp=\{y:B(y,e_1)=0\}$, so $V=Ke_1\oplus e_1^\perp$, an orthogonal sum, and $f\cong\langle a\rangle\oplus f|_{e_1^\perp}$.
Induct on $n$.

(b) Let $x\ne0$ with $f(x)=0$. Since $f$ is non-degenerate there is $y$ with $B(x,y)=1$. Then $f(y+tx)=f(y)+2t$, and for
$t=-f(y)/2$ we get $y'=y+tx$ with $f(y')=0$, $B(x,y')=1$. On the plane $H=Kx+Ky'$, $f(ux+vy')=2uv$, which is non-degenerate, so
$V=H\oplus H^\perp$ as in (a). The form $2uv$ represents $1$ (at $u=\frac12$, $v=1$), so by (a) it is $\langle1,c\rangle$ for
some $c$; its Gram matrix $\bigl(\begin{smallmatrix}0&1\\1&0\end{smallmatrix}\bigr)$ has determinant $-1$, so $c=-1$ up to
squares and $H\cong\langle1,-1\rangle$. Finally $2uv$ takes the value $a$ at $(u,v)=(a/2,1)$.

(c) If $f(x)=a$, then $f(x)-a\cdot1^2=0$. Conversely if $f(x)-at^2=0$ with $(x,t)\ne0$: if $t\ne0$, $f(x/t)=a$; if $t=0$, $f$ is
isotropic and represents $a$ by (b).
\end{proof}

\section{The Hilbert symbol}

For $a,b\in\Q_v^\times$, put $(a,b)_v=1$ if $z^2=ax^2+by^2$ has a solution $(x,y,z)\ne0$ in $\Q_v^3$, and $(a,b)_v=-1$ otherwise.
It depends only on the classes of $a,b$ modulo squares. For $v=\infty$, clearly $(a,b)_\infty=-1$ iff $a<0$ and $b<0$.

\begin{proposition}\label{prop:hilbertbasic}
For $a,a',b\in\Q_v^\times$:
\begin{enumerate}[label=(\alph*)]
\item $(a,b)=(b,a)$, $(a,c^2)=1$, $(a,-a)=1$ and $(a,1-a)=1$ (for $a\ne1$).
\item If $b$ is not a square, $(a,b)=1$ iff $a$ is a norm from $\Q_v(\sqrt b)$, that is, $a=z^2-by^2$ with $z,y\in\Q_v$.
\item If $(a,b)=1$, then $(aa',b)=(a',b)$.
\item $(a,b)=(a,-ab)=(a,(1-a)b)$.
\end{enumerate}
\end{proposition}
\begin{proof}
(a) Symmetry and $(a,c^2)=1$ ($x=0$, $y=1$, $z=c$) are clear; $(1,1)$ on $z^2=ax^2-ay^2$ and on $z^2=ax^2+(1-a)y^2$.
(b) If $z^2=ax^2+by^2$ with $x=0$, then $z^2=by^2$ forces $y=z=0$; so $x\ne0$ and $a=N\bigl(\frac zx+\frac yx\sqrt b\bigr)$.
Conversely $a=N(z+y\sqrt b)$ gives the solution $(1,y,z)$. (c) If $b$ is a square both sides are $1$. Otherwise the norms
form a subgroup $N$ of $\Q_v^\times$, and $a\in N$ implies ($a'\in N\iff aa'\in N$). (d) Apply (c), with the roles of the two
arguments exchanged, to $(a,-a)=1$ and $(a,1-a)=1$.
\end{proof}

\begin{theorem}[the formulas]\label{thm:hilbertformula}
Write $a=p^\alpha u$, $b=p^\beta v$ with $u,v$ units of $\Z_p$.
\begin{enumerate}[label=(\alph*)]
\item If $p$ is odd, $(a,b)_p=(-1)^{\alpha\beta\frac{p-1}2}\bigl(\frac up\bigr)^\beta\bigl(\frac vp\bigr)^\alpha$.
\item If $p=2$, $(a,b)_2=(-1)^{\epsilon(u)\epsilon(v)+\alpha\omega(v)+\beta\omega(u)}$, where $\epsilon(u)=\frac{u-1}2$ and
$\omega(u)=\frac{u^2-1}8$ modulo $2$.
\end{enumerate}
\end{theorem}

Here $\bigl(\frac up\bigr)$ means the Legendre symbol of $u\bmod p$, and $\epsilon(u)$, $\omega(u)$ depend only on $u\bmod 8$:
$\epsilon(u)=0$ iff $u\equiv1\pmod4$, and $\omega(u)=0$ iff $u\equiv\pm1\pmod8$. Both are additive: $\epsilon(uv)=\epsilon(u)+
\epsilon(v)$ and $\omega(uv)=\omega(u)+\omega(v)$ modulo $2$ (check the four classes modulo $8$). Figure~\ref{fig:hilbert2} shows
the table for $p=2$.

\begin{proof}
Both sides depend only on $\alpha,\beta$ modulo $2$ (for the left side, replace $a$ by $a/p^2$), and both are symmetric, so
there are three cases.

\emph{(a), $\alpha=\beta=0$.} We show $(u,v)_p=1$. The sets $\{ux^2\}$ and $\{1-vy^2\}$, for $x,y\in\mathbb F_p$, each have
$\frac{p+1}2$ elements, so they meet: $ux_0^2+vy_0^2\equiv1\pmod p$, with $x_0$ or $y_0$ non-zero, say $x_0$. Hensel's lemma
applied to $X\mapsto uX^2+vy_0^2-1$ at $x_0$ (derivative $2ux_0$, a unit) gives $x\in\Z_p$ with $ux^2+vy_0^2=1$.

\emph{(a), $\alpha=1$, $\beta=0$.} We show $(pu,v)_p=\bigl(\frac vp\bigr)$. If $v$ is a square modulo $p$, it is a square in
$\Q_p$ and the symbol is $1$. Conversely let $z^2=pux^2+vy^2$ with $(x,y,z)$ primitive (in $\Z_p$, not all divisible by $p$).
If $p\mid y$ and $p\mid z$, then $p^2\mid pux^2$, so $p\mid x$: impossible. So $z^2\equiv vy^2\pmod p$ with $y\not\equiv0$, and
$v$ is a square modulo $p$.

\emph{(a), $\alpha=\beta=1$.} By Proposition~\ref{prop:hilbertbasic}(d), $(pu,pv)=(pu,-p^2uv)=(pu,-uv)=\bigl(\frac{-uv}p\bigr)=
(-1)^{\frac{p-1}2}\bigl(\frac up\bigr)\bigl(\frac vp\bigr)$.

\emph{(b), $\alpha=\beta=0$.} If $u\equiv1\pmod8$, $u$ is a square. If $u\equiv5\pmod8$, then $u+4v\equiv1\pmod8$ is a square
$w^2$, and $(x,y,z)=(1,2,w)$ solves $z^2=ux^2+vy^2$. So $(u,v)=1$ if $u\equiv1\pmod4$, and likewise if $v\equiv1\pmod4$.
If $u\equiv v\equiv3\pmod4$ and $z^2=ux^2+vy^2$ with $(x,y,z)$ primitive, then $z^2+x^2+y^2\equiv0\pmod4$; squares are
$0,1$ modulo $4$, so $x,y,z$ are all even, a contradiction. So $(u,v)=(-1)^{\epsilon(u)\epsilon(v)}$.

\emph{(b), $\alpha=1$, $\beta=0$.} First, $(2,v)=(-1)^{\omega(v)}$. If $v\equiv1\pmod 8$ it is a square; if $v\equiv-1$, then
$2+v\equiv1$ is a square $w^2$ and $(1,1,w)$ is a solution. Conversely, if $z^2=2x^2+vy^2$ is primitive, then $y,z$ are odd
(if $y$ is even, then $z$ is even, then $x$ is even), so $z^2\equiv y^2\equiv1\pmod8$ and $1\equiv2x^2+v\pmod 8$; as $2x^2\equiv
0$ or $2$, $v\equiv\pm1\pmod8$. Next, $(2u,v)=(2,v)(u,v)$: by Proposition~\ref{prop:hilbertbasic}(c) this holds if $(2,v)=1$ or
$(u,v)=1$; in the remaining case $v\equiv3\pmod8$ and $u\equiv3\pmod4$, and then $2u+v\equiv6+3$ or $14+3$, which is
$\equiv1\pmod 8$, so $(1,1,\sqrt{2u+v})$ shows $(2u,v)=1=(2,v)(u,v)$. Hence $(2u,v)=(-1)^{\epsilon(u)\epsilon(v)+\omega(v)}$.

\emph{(b), $\alpha=\beta=1$.} $(2u,2v)=(2u,-4uv)=(2u,-uv)$, and the previous case with $-uv$ in place of $v$ gives
$(-1)^{\epsilon(u)\epsilon(-uv)+\omega(-uv)}$. Now $\epsilon(-uv)=1+\epsilon(u)+\epsilon(v)$ and $\omega(-uv)=\omega(u)+\omega(v)$
($\omega(-1)=0$), and $\epsilon(u)(1+\epsilon(u))=0$; so the exponent is $\epsilon(u)\epsilon(v)+\omega(u)+\omega(v)$, as claimed.
\end{proof}

\begin{figure}[ht]
\centering
\fig{ch07-hilbert2}
\caption{The Hilbert symbol $(a,b)_2$ on the eight square classes of $\Q_2$. Each row has four $+$ and four $-$, except the
row of $1$: the symbol is a non-degenerate symmetric bilinear form on $\Q_2^\times/(\Q_2^\times)^2\cong(\Z/2)^3$.}
\label{fig:hilbert2}
\end{figure}

\begin{corollary}\label{cor:hilbertbilinear}
$(a,b)_v$ is bimultiplicative: $(aa',b)_v=(a,b)_v(a',b)_v$. It is non-degenerate: if $(a,b)_v=1$ for all $b$, then $a$ is a
square. Consequently, if $a$ is not a square, exactly half of the square classes $b$ have $(a,b)_v=1$.
\end{corollary}
\begin{proof}
For $v=\infty$ it is clear. For $p$, the formulas are visibly multiplicative in each argument, since the Legendre symbol,
$\epsilon$ and $\omega$ are. Non-degeneracy: for $p$ odd, if $a=p^\alpha u$ is not a square then either $\alpha$ is odd and
$b=v$ a non-residue unit gives $-1$, or $\alpha$ is even, $u$ a non-residue, and $b=p$ gives $-1$. For $p=2$ read
Figure~\ref{fig:hilbert2}: every row but the first contains a $-1$. Finally $b\mapsto(a,b)_v$ is then a non-trivial
homomorphism to $\{\pm1\}$, whose kernel has index $2$.
\end{proof}

\section{Hilbert's product formula}

\begin{theorem}[Hilbert]\label{thm:hilbertproduct}
For $a,b\in\Q^\times$, $(a,b)_v=1$ for all but finitely many places $v$, and $\displaystyle\prod_v(a,b)_v=1$.
\end{theorem}
\begin{proof}
For $p$ odd not dividing the numerators and denominators of $a,b$, both are units at $p$ and $(a,b)_p=1$. By
bimultiplicativity it suffices to treat $a,b\in\{-1,2\}\cup\{\text{odd primes}\}$. Write $\ell,\ell'$ for odd primes.
\begin{itemize}
\item $(-1,-1)$: $-1$ at $\infty$ and at $2$ ($\epsilon(-1)=1$), $1$ elsewhere.
\item $(-1,2)$: $1$ everywhere ($\omega(-1)=0$ at $2$).
\item $(2,2)=(2,-1)$: $1$ everywhere.
\item $(-1,\ell)$: $\bigl(\frac{-1}\ell\bigr)$ at $\ell$, $(-1)^{\epsilon(\ell)}$ at $2$, which agree by the first supplementary law.
\item $(2,\ell)$: $\bigl(\frac2\ell\bigr)$ at $\ell$, $(-1)^{\omega(\ell)}$ at $2$, which agree by the second supplementary law.
\item $(\ell,\ell)=(\ell,-1)$, done above.
\item $(\ell,\ell')$ with $\ell\ne\ell'$: $\bigl(\frac{\ell'}\ell\bigr)$ at $\ell$, $\bigl(\frac\ell{\ell'}\bigr)$ at $\ell'$,
$(-1)^{\epsilon(\ell)\epsilon(\ell')}$ at $2$. Their product is $1$: this is quadratic reciprocity.
\end{itemize}
\end{proof}

The proof shows that the product formula \emph{is} quadratic reciprocity, with its two supplements, rewritten so that the
prime $2$ and the real place play the same role as the odd primes (Exercise~\ref{ex:prodQR}). Here is the hook of the chapter.
For $f=\langle2,7,-3\rangle$, dividing by $3$ shows that $f$ is isotropic over $\Q_v$ iff $z^2=\frac23x^2+\frac73y^2$ has a
solution, iff $(6,21)_v=1$. By Theorem~\ref{thm:hilbertformula}: $(6,21)_\infty=1$; at $3$, $(3\cdot2,3\cdot7)_3=(-1)^{1}
\bigl(\frac23\bigr)\bigl(\frac73\bigr)=(-1)(-1)(1)=1$; at $7$, $(6,7)_7=\bigl(\frac67\bigr)=-1$; at $2$, $(2\cdot3,21)_2=
(-1)^{\epsilon(3)\epsilon(21)+\omega(21)}=(-1)^{0+1}=-1$; at all other $p$ both are units. So $f$ fails exactly at $2$ and $7$.
The product formula says that the number of failures is even: \emph{a single failure is impossible}, and once we know the
answer at all places but one, we know it at the last.

\section{Quadratic forms over $\Q_p$}

Let $f\cong\langle a_1,\dots,a_n\rangle$ over $\Q_v$, and define the \emph{Hasse invariant}
\[
\varepsilon(f)=\prod_{i<j}(a_i,a_j)_v.
\]

\begin{proposition}\label{prop:hasseinv}
$\varepsilon(f)$ does not depend on the diagonalisation.
\end{proposition}
\begin{proof}
By bimultiplicativity, $\varepsilon(\langle a_1\rangle\oplus g)=(a_1,d(g))\,\varepsilon(g)$ for diagonal $g$. We argue by induction on
$n$; $n=1$ is trivial.

\emph{$n=2$.} Suppose $\langle a,b\rangle\cong\langle a',b'\rangle$. Then $ab=a'b'$ modulo squares, and $\langle a,b\rangle$
represents $a'$, so $\langle a,b,-a'\rangle$ is isotropic, that is, $z^2=\frac a{a'}x^2+\frac b{a'}y^2$ has a solution, so
$(aa',ba')=1$. Expanding, $(a,b)(a,a')(a',b)(a',a')=1$, so $(a,b)=(a',ab)(a',a')=(a',ab)(a',-1)=(a',-ab)$, using $(a',a')=(a',-1)$
(Proposition~\ref{prop:hilbertbasic}(d)). And $(a',b')=(a',a'ab)=(a',-1)(a',ab)$ by the same identity, which is the same number.

\emph{$n\ge3$.} Call two orthogonal bases \emph{contiguous} if they share a vector. If $(e_i)$ and $(e'_i)$ share, say,
$e_1=e'_1$, then both diagonalise the restriction of $f$ to $e_1^\perp$, with the same discriminant $d(f)/f(e_1)$, and by
induction the same Hasse invariant; so $\varepsilon$ is the same for both bases. It remains to connect any two orthogonal bases
$(e_i)$, $(e'_i)$ by a chain of contiguous ones. Write $e'_1=\sum_ix_ie_i$, and induct on the number $k$ of non-zero $x_i$. If
$k=1$, $e'_1$ is a multiple of $e_1$, and the basis $(e'_1,e_2,\dots,e_n)$ is contiguous to $(e_i)$ (through $e_2$) and to
$(e'_i)$ (through $e'_1$). If $k\ge2$, the numbers $t_i=x_i^2f(e_i)$ ($x_i\ne0$) have non-zero sum $f(e'_1)$; for $k=2$ this
says $t_1+t_2\ne0$, and for $k\ge3$ some pair has $t_i+t_j\ne0$ (if $t_1+t_2=t_1+t_3=0$ then $t_2+t_3=2t_2\ne0$). Renumber so
that $f(x_1e_1+x_2e_2)=t_1+t_2\ne0$, and complete $\varepsilon_1=x_1e_1+x_2e_2$ to an orthogonal basis $(\varepsilon_1,\varepsilon_2)$
of $Ke_1+Ke_2$. The basis $(\varepsilon_1,\varepsilon_2,e_3,\dots,e_n)$ is contiguous to $(e_i)$ through $e_3$, and in it $e'_1$ has
$k-1$ non-zero coordinates; by induction it is chain-connected to $(e'_i)$.
\end{proof}

\begin{theorem}[forms over $\Q_v$]\label{thm:localforms}
Let $f$ be a non-degenerate form of rank $n$ over $\Q_p$, $d=d(f)$, $\varepsilon=\varepsilon(f)$.
\begin{enumerate}[label=(\alph*)]
\item $f$ is isotropic if and only if: $n=2$ and $d=-1$; or $n=3$ and $(-1,-d)=\varepsilon$; or $n=4$ and ($d\ne1$ or
$\varepsilon=(-1,-1)$); or $n\ge5$.
\item $f$ represents $a\in\Q_p^\times$ if and only if: $n=1$ and $a=d$; or $n=2$ and $(a,-d)=\varepsilon$; or $n=3$ and ($a\ne-d$ or
$(-1,-d)=\varepsilon$); or $n\ge4$. (Classes modulo squares throughout.)
\item Two forms over $\Q_p$ are equivalent if and only if they have the same rank, discriminant and Hasse invariant.
\end{enumerate}
\end{theorem}
\begin{proof}
(a) $n=2$: $\langle a,b\rangle$ is isotropic iff $-b/a=-ab/a^2$ is a square, iff $d=-1$. $n=3$: $\langle a,b,c\rangle$ is isotropic iff
$z^2=-\frac acx^2-\frac bcy^2$ has a solution, iff $(-ac,-bc)=1$. Expanding with bimultiplicativity and $(c,c)=(c,-1)$,
\[
(-ac,-bc)=(-1,-1)(-1,a)(-1,b)(-1,c)\,(a,b)(a,c)(b,c)=(-1,-abc)\,\varepsilon=(-1,-d)\,\varepsilon,
\]
which is $1$ iff $(-1,-d)=\varepsilon$.

For (b), $f$ represents $a$ iff $f\oplus\langle-a\rangle$ is isotropic (Proposition~\ref{prop:diag}(c)), a form with
discriminant $-ad$ and Hasse invariant $\varepsilon\,(-a,d)$. For $n=2$, the criterion of (a) for $n=3$ reads $(-1,ad)=
\varepsilon(-a,d)$, that is $\varepsilon=(-1,ad)(-a,d)=(-1,a)(-1,d)(-1,d)(a,d)=(a,-d)$, as claimed. For $n=1$ it is clear.

(a) $n=4$: write $f=\langle a_1,a_2\rangle\oplus\langle a_3,a_4\rangle$. Then $f$ is isotropic iff some $x\in\Q_p^\times$ is
represented both by $g=\langle a_1,a_2\rangle$ and by $-h=\langle-a_3,-a_4\rangle$ (if $g(y)+h(y')=0$ with $(y,y')\ne0$ and the
common value is $0$, one of $g,-h$ is isotropic, hence represents everything, and the other represents some $x\ne0$). By (b)
for $n=2$, this means $(x,-a_1a_2)=(a_1,a_2)$ and $(x,-a_3a_4)=(-a_3,-a_4)$. If $-a_1a_2$ is a square, $g$ is isotropic and we are
done; likewise for $-a_3a_4$. Otherwise each condition holds for exactly half the classes $x$ (Corollary~\ref{cor:hilbertbilinear}),
the sets being cosets of the kernels of the characters $x\mapsto(x,-a_1a_2)$ and $x\mapsto(x,-a_3a_4)$. If these characters are
different, the two cosets meet (each character takes both values on the kernel of the other, since a group is not the union of
two proper subgroups). If they are equal, that is $a_1a_2=a_3a_4$, i.e.\ $d=1$, the cosets meet iff $(a_1,a_2)=(-a_3,-a_4)$.
When $d=1$ a short computation with $a_3a_4=a_1a_2$ gives $\varepsilon=(a_1,a_2)(a_3,a_4)(a_1a_2,-1)$ and $(-a_3,-a_4)=(-1,-1)
(-1,a_1a_2)(a_3,a_4)$, so $(a_1,a_2)(-a_3,-a_4)=\varepsilon\,(-1,-1)$. Hence: isotropic iff $d\ne1$ or $\varepsilon=(-1,-1)$.

$n\ge5$: write $f=g\oplus h$ with $g$ of rank $2$ and $h$ of rank $n-2\ge3$. By (b) for $n=3$ (and its trivial extension to rank
$\ge3$ by adding variables), $-h$ represents all classes except at most one; $g$ represents at least two classes (half of the at
least four classes, or all of them). So some $x$ is represented by both, and $f$ is isotropic.

(b) for $n=3$ and $n\ge4$ follows from (a) for $n=4$ and $n\ge5$ applied to $f\oplus\langle-a\rangle$: for $n=3$ the form has
discriminant $-ad$, which is $1$ iff $a=-d$, and then its Hasse invariant is $\varepsilon(-a,d)=\varepsilon(d,d)=\varepsilon(-1,d)$, and the
condition $\varepsilon(-1,d)=(-1,-1)$ is $\varepsilon=(-1,-d)$.

(c) Equivalent forms have the same invariants (Proposition~\ref{prop:hasseinv}). Conversely, if $f,g$ have the same $n,d,\varepsilon$,
then by (b) they represent the same classes; pick $a$ represented by both, write $f\cong\langle a\rangle\oplus f'$,
$g\cong\langle a\rangle\oplus g'$ (Proposition~\ref{prop:diag}(a)). Then $d(f')=d(g')=d/a$ and $\varepsilon(f')=(a,d/a)\varepsilon(f)=
\varepsilon(g')$, so $f'\cong g'$ by induction on $n$, and $f\cong g$.
\end{proof}

\begin{think}
By (a), every form in five or more variables over $\Q_p$ has a non-trivial zero, while four can fail: $\langle1,1,1,1\rangle$ has
$d=1$ and $\varepsilon=1\ne(-1,-1)_2$, so $x_1^2+x_2^2+x_3^2+x_4^2=0$ has no non-trivial solution in $\Q_2$ (check it directly:
$-1\equiv7\pmod 8$ is not a sum of three squares in $\Q_2$). Compare with $\R$, where $x_1^2+\dots+x_n^2$ is anisotropic in every
dimension. Why can a $p$-adic field not carry a ``positive definite'' form in five variables? (Hint: the residue field is
finite, and Chevalley's theorem says that a form of degree $2$ in $3$ variables over $\mathbb F_p$ has a non-trivial zero.)
\end{think}

\section{The Hasse--Minkowski theorem}

We need two lemmas about $\Q$ inside all its completions.

\begin{lemma}[approximation]\label{lem:approx}
Let $S$ be a finite set of places of $\Q$, and $x_v\in\Q_v$ for $v\in S$. For every $\delta>0$ there is $x\in\Q$ with
$|x-x_v|_v<\delta$ for all $v\in S$.
\end{lemma}
\begin{proof}
Choose a positive integer $N$ with $Nx_p\in\Z_p$ for every prime $p\in S$. If $x'\in\Q$ satisfies $|x'-Nx_v|_v<\delta'$ for all
$v\in S$, then $x=x'/N$ satisfies $|x-x_v|_v=|N|_v^{-1}|x'-Nx_v|_v<\delta'/|N|_v$, which is at most $\delta$ for
$\delta'=\delta\min_{v\in S}|N|_v$. So we may assume $x_p\in\Z_p$ for the primes $p\in S$. Choose $m$ with $p^{-m}<\delta$. By the
Chinese remainder theorem there is $z\in\Z$ with $z\equiv x_p\pmod{p^m}$ for all primes $p\in S$; let $M=\prod p^m$. For a prime
$q\notin S$ and $c\in\Z$, the number $z+Mc/q$ is still within $p^{-m}$ of $x_p$ at each $p\in S$, and these numbers are spaced
$M/q$ apart on the real line; choose $q>M/\delta$ and then $c$ to get within $\delta$ of $x_\infty$ (if $\infty\in S$).
\end{proof}

\begin{theorem}[existence of a rational number with given Hilbert symbols]\label{thm:existence}
Let $a_1,\dots,a_m\in\Q^\times$ and $\epsilon_{i,v}=\pm1$ for $1\le i\le m$ and all places $v$. There is $x\in\Q^\times$ with
$(a_i,x)_v=\epsilon_{i,v}$ for all $i$ and $v$ if and only if
\begin{enumerate}[label=(\arabic*)]
\item almost all $\epsilon_{i,v}$ equal $1$;
\item $\prod_v\epsilon_{i,v}=1$ for each $i$;
\item for each $v$ there is $x_v\in\Q_v^\times$ with $(a_i,x_v)_v=\epsilon_{i,v}$ for all $i$.
\end{enumerate}
\end{theorem}
\begin{proof}
The conditions are necessary by Theorem~\ref{thm:hilbertproduct}. Conversely, we may take the $a_i$ to be integers. Let $S$ be
the set made of $\infty$, $2$, the primes dividing some $a_i$, and the places where some $\epsilon_{i,v}=-1$; it is finite by
(1). By Lemma~\ref{lem:approx}, choose $x'\in\Q^\times$ so close to $x_v$ for every $v\in S$ that $x'/x_v$ is a square in $\Q_v$ (the
squares form an open subgroup containing $1$); then $(a_i,x')_v=\epsilon_{i,v}$ for $v\in S$. Multiplying $x'$ by the square of
its denominator, we may write $x'=t\,u$ with $t=\pm\prod_{p\in S}p^{\alpha_p}$ and $u$ a positive integer prime to every prime of
$S$. Let $M=8\prod_{p\in S,\,p\ne2,\infty}p$. By Dirichlet's theorem there is a prime $\ell\notin S$ with $\ell\equiv u\pmod M$.
Put $x=t\ell$.

For $v\in S$: $\ell/u$ is positive, $\equiv1\pmod8$ and $\equiv1$ modulo each odd $p\in S$, so it is a square in every $\Q_v$,
$v\in S$; hence $(a_i,x)_v=(a_i,x')_v=\epsilon_{i,v}$. For $v\notin S\cup\{\ell\}$: $v$ is odd and $a_i,x$ are $v$-adic units, so
$(a_i,x)_v=1=\epsilon_{i,v}$. For $v=\ell$: by the product formula and (2),
$(a_i,x)_\ell=\prod_{v\ne\ell}(a_i,x)_v=\prod_{v\ne\ell}\epsilon_{i,v}=\epsilon_{i,\ell}$.
\end{proof}

\begin{theorem}[Hasse--Minkowski]\label{thm:hm}
A non-degenerate quadratic form with rational coefficients is isotropic over $\Q$ if and only if it is isotropic over $\R$
and over $\Q_p$ for every prime $p$.
\end{theorem}
\begin{proof}
Only ``if'' needs proof. Diagonalise, $f=\langle a_1,\dots,a_n\rangle$; multiplying by a constant, $a_1=1$.

\emph{$n=2$.} $f$ is isotropic iff $-a_2$ is a square. A rational number that is a square in every $\Q_p$ and in $\R$ has even
valuation at every $p$ and is positive, so it is a square.

\emph{$n=3$.} We may assume $f=z^2-ax^2-by^2$ with $a,b$ squarefree integers, and argue by induction on $|a|+|b|$, assuming
$|a|\le|b|$. If $|b|=1$, then $(a,b)_\infty=1$ excludes $a=b=-1$, so $a=1$ or $b=1$, and $f$ is isotropic. Let $|b|\ge2$ and $p\mid
b$. We claim $a$ is a square modulo $p$. If $p\mid a$, then $a\equiv0=0^2$. If $p=2$, every integer is a square modulo $2$.
If $p$ is odd and $p\nmid a$, then $a$ is a unit and $v_p(b)=1$, so $1=(a,b)_p=\bigl(\frac ap\bigr)$ by
Theorem~\ref{thm:hilbertformula}. So $a$ is a square modulo every $p\mid b$, hence modulo $b$ (Chinese remainder theorem,
$b$ squarefree): $a\equiv t^2\pmod b$ with $|t|\le|b|/2$. Write $t^2-a=bb'$. If $b'=0$, $a$ is a square and $f$ is isotropic.
Otherwise $bb'=N(t+\sqrt a)$ is a norm from $\Q(\sqrt a)$ (and from each $\Q_v(\sqrt a)$), so $b$ is a norm iff $b'$ is: $f$ is
isotropic iff $f'=z^2-ax^2-b'y^2$ is, over $\Q$ and over each $\Q_v$ (Proposition~\ref{prop:hilbertbasic}(b),(c)). Now
$|b'|=\frac{|t^2-a|}{|b|}\le\frac{|b|}4+1<|b|$. Write $b'=b''w^2$ with $b''$ squarefree; then $f'$ is equivalent to
$z^2-ax^2-b''y^2$ with $|a|+|b''|<|a|+|b|$, isotropic everywhere locally, hence isotropic over $\Q$ by induction.

\emph{$n=4$.} Write $f=\langle a_1,a_2\rangle\oplus\langle-a_3,-a_4\rangle=g-h$. For each $v$, as $f$ is isotropic over $\Q_v$, there
is $x_v\in\Q_v^\times$ represented by both $g$ and $h$ (as in the proof of Theorem~\ref{thm:localforms}(a)), that is,
$(x_v,-a_1a_2)_v=(a_1,a_2)_v$ and $(x_v,-a_3a_4)_v=(a_3,a_4)_v$. The numbers $\epsilon_{1,v}=(a_1,a_2)_v$, $\epsilon_{2,v}=(a_3,a_4)_v$
satisfy (1)--(3) of Theorem~\ref{thm:existence} with $-a_1a_2,-a_3a_4$, so there is $x\in\Q^\times$ with $(x,-a_1a_2)_v=(a_1,a_2)_v$
and $(x,-a_3a_4)_v=(a_3,a_4)_v$ for all $v$. Then $\langle a_1,a_2,-x\rangle$ is isotropic over every $\Q_v$ (Theorem~\ref{thm:localforms}(b),
or its obvious analogue over $\R$), hence over $\Q$ by the case $n=3$: $g$ represents $x$ over $\Q$. Likewise $h$ represents $x$,
and $f=g-h$ is isotropic.

\emph{$n\ge5$.} Write $f=g-h$ with $g=\langle a_1,a_2\rangle$ and $h$ of rank $n-2\ge3$. Let $S$ consist of $\infty$, $2$ and the
primes $p$ at which some $a_i$ is not a unit. For $v\in S$, pick $x_v\ne0$ represented over $\Q_v$ by $g$ and $h$, say $x_v=g(y_v)$.
By Lemma~\ref{lem:approx}, applied to both coordinates, choose $y\in\Q^2$ so close to each $y_v$ that $x=g(y)\ne0$ satisfies
$x/x_v\in\Q_v^{\times2}$ for $v\in S$. Consider $f_1=\langle x\rangle-h$, of rank $n-1\ge4$. For $v\in S$ it is isotropic, since $h$
represents $x_v$, hence $x$. For $v\notin S$, $h$ has unit coefficients at the odd prime $v$, so any three of them form an
isotropic form ($d$ is a unit and $\varepsilon=1$, and $(-1,-d)_v=1$; Theorem~\ref{thm:localforms}(a)); so $f_1$ is isotropic.
By induction (the cases $4,5,\dots,n-1$), $f_1$ is isotropic over $\Q$, so $h$ represents $x$ over $\Q$ (Proposition~\ref{prop:diag}(c)),
and so does $g$. Hence $f$ is isotropic.
\end{proof}

\begin{corollary}\label{cor:hmrep}
A form with rational coefficients represents $a\in\Q^\times$ over $\Q$ iff it represents $a$ over $\R$ and over every $\Q_p$.
\end{corollary}
\begin{proof}
Apply the theorem to $f\oplus\langle-a\rangle$ (Proposition~\ref{prop:diag}(c)).
\end{proof}

\begin{corollary}[Meyer, 1884]\label{cor:meyer}
An indefinite form with rational coefficients in at least five variables has a non-trivial rational zero.
\end{corollary}
\begin{proof}
It is isotropic over $\R$ (indefinite) and over every $\Q_p$ (Theorem~\ref{thm:localforms}(a), $n\ge5$).
\end{proof}

\section{Legendre's theorem, and three squares}

\begin{theorem}[Legendre, 1785]\label{thm:legendre3}
Let $a,b,c$ be non-zero squarefree integers, pairwise coprime and not all of the same sign. Then $ax^2+by^2+cz^2=0$ has a
non-trivial integer solution if and only if $-bc$ is a square modulo $|a|$, $-ca$ modulo $|b|$, and $-ab$ modulo $|c|$.
\end{theorem}
\begin{proof}
The form is isotropic over $\Q_v$ iff $(-ac,-bc)_v=1$ (proof of Theorem~\ref{thm:localforms}(a)), and a rational zero can be scaled to
an integer one. At $\infty$ the signs are mixed, so the symbol is $1$. At an odd $p\nmid abc$, both arguments are units: $1$. At
an odd $p\mid a$, write $a=pa'$; then $-ac=p(-a'c)$ and $-bc$ is a unit, and Theorem~\ref{thm:hilbertformula} gives $(-ac,-bc)_p=
\bigl(\frac{-bc}p\bigr)$; similarly at $p\mid b$ and $p\mid c$. By the Chinese remainder theorem the stated conditions are exactly
that these symbols are $1$ at all odd $p\mid abc$ (modulo $2$ every number is a square). Then all symbols except possibly
$(-ac,-bc)_2$ are $1$, and the product formula forces that one to be $1$ too. Conclude with Hasse--Minkowski.
\end{proof}

\begin{theorem}[Gauss, Legendre]\label{thm:threesq}
A positive integer is a sum of three squares of integers if and only if it is not of the form $4^a(8b+7)$.
\end{theorem}

The ``only if'' part is Tome~1, Exercise~7.1 (squares modulo $8$). For the converse we first get rational solutions and then
integer ones.

\begin{lemma}[Aubry, Davenport--Cassels]\label{lem:dc}
Let $f(x)=\sum a_{ij}x_ix_j$ be a positive definite form with $a_{ij}=a_{ji}\in\Z$, such that for every $x\in\Q^n$ there is $y\in
\Z^n$ with $f(x-y)<1$. If an integer $m$ is represented by $f$ over $\Q$, it is represented over $\Z$.
\end{lemma}
\begin{proof}
Let $B$ be the bilinear form, with $B(y,z)\in\Z$ for $y,z\in\Z^n$. Write a rational solution as $x=y/t$ with $y\in\Z^n$, $t\ge1$
an integer, $f(y)=mt^2$, and $t$ minimal. Suppose $t\ge2$; then $x\notin\Z^n$. Choose $z\in\Z^n$ with $0<f(x-z)<1$, and put
\[
a=f(z)-m,\qquad b=2\bigl(mt-B(y,z)\bigr),\qquad t'=at+b,\qquad y'=ay+bz.
\]
These are integers. A direct computation gives $tf(x-z)=tm-2B(y,z)+tf(z)=t'$, so $0<t'<t$. And
\[
f(y')=a^2mt^2+2abB(y,z)+b^2f(z),\qquad mt'^2=a^2mt^2+2abmt+b^2m,
\]
whose difference is $ab\bigl(2B(y,z)-2mt\bigr)+b^2(f(z)-m)=-ab^2+ab^2=0$. So $y'/t'$ is a rational solution with smaller
denominator, contradicting the minimality of $t$.
\end{proof}

\begin{proof}[Proof of Theorem~\ref{thm:threesq}]
Let $n=4^am$ with $m\not\equiv7\pmod8$ (and $4\nmid m$ or not, it does not matter). By Corollary~\ref{cor:hmrep} we check that
$\langle1,1,1,-n\rangle$ is isotropic over every $\Q_v$. Over $\R$ it is indefinite. Over $\Q_p$, $p$ odd, the unit form
$\langle1,1,1\rangle$ is already isotropic ($d=1$, $\varepsilon=1=(-1,-1)_p$). Over $\Q_2$, by Theorem~\ref{thm:localforms}(a) it
suffices that the discriminant $-n$ is not a square; and $-n=4^a(-m)$ is a square in $\Q_2$ iff $-m\equiv1\pmod8$, i.e.\ iff
$m\equiv7\pmod8$, which we excluded. So $n$ is a sum of three rational squares. For $f=x^2+y^2+z^2$ and $x\in\Q^3$, rounding each
coordinate to the nearest integer gives $y$ with $f(x-y)\le\frac34<1$, so Lemma~\ref{lem:dc} applies.
\end{proof}

\begin{think}
Ernst Selmer showed in 1951 that the cubic equation $3x^3+4y^3+5z^3=0$ has non-zero solutions in $\R$ and in every $\Q_p$, but
none in $\Q$ (Exercise~\ref{ex:selmer} proves the local part at most primes). Where does the proof of Hasse--Minkowski use that the
degree is $2$? (Diagonalisation, the Hilbert symbol and its product formula, and above all the fact that a quadric with one
rational point is parametrised by lines through it. For cubics the obstruction is measured by the Tate--Shafarevich group of an
elliptic curve; Chapter~12.)
\end{think}

\section{Exercises}

\begin{exercise}[$\star$]\label{ex:fourrat}
Show that every positive rational number is a sum of four squares of rationals, and that $r>0$ is a sum of three squares of
rationals if and only if $-r$ is not a square in $\Q_2$. Is $7$ a sum of three rational squares? Is $\frac7{4}$? Is $28$?
\end{exercise}

\begin{exercise}[$\star$]\label{ex:fivesixths}
Show that the proportion of $n\le x$ that are sums of three squares tends to $\frac56$.
\end{exercise}

\begin{exercise}[$\star\star$]\label{ex:eureka}
\emph{Gauss's Eureka.} Deduce from Theorem~\ref{thm:threesq} that every positive integer is a sum of three triangular numbers
$\frac{k(k+1)}2$, $k\ge0$.
\end{exercise}

\begin{exercise}[$\star\star$]\label{ex:legendre13}
Does $13x^2+7y^2=z^2$ have a non-trivial rational solution? At which places does it fail, and how does the product formula show
up?
\end{exercise}

\begin{exercise}[$\star\star$]\label{ex:prodQR}
Show conversely that the product formula, for $a,b\in\{-1,2,\ell,\ell'\}$, implies quadratic reciprocity and its two supplements,
given the local formulas of Theorem~\ref{thm:hilbertformula}.
\end{exercise}

\begin{exercise}[$\star\star$]\label{ex:x2y3z}
Show that $x^2+2y^2+3z^2$ represents a positive integer $n$ over $\Q$ if and only if $n$ is not of the form $4^a(16b+10)$.
(Ramanujan's form; it is \emph{regular}: the same holds over $\Z$.)
\end{exercise}

\begin{exercise}[$\star\star\star$]\label{ex:selmer}
\emph{Selmer's cubic, locally.} Show that $3x^3+4y^3+5z^3=0$ has a non-zero solution in $\R$, in $\Q_2$, $\Q_3$, $\Q_5$, and in $\Q_p$
for every $p\equiv2\pmod3$.
\end{exercise}

\begin{beauty}
Two numbers, $\pm1$, at each place, and a product formula tying them together: this is all it takes to decide every quadratic
equation over $\Q$. The local--global principle was the first great success of the idea that a global object is governed by
its completions, and it became the model for everything that followed: class field theory is a local--global principle for
abelian extensions, the ad\`eles make it into a space, and the failures of the principle, for cubics and beyond, became a
subject of their own (Brauer--Manin, Tate--Shafarevich). The product formula itself, once read as reciprocity, is the first
case of Artin's reciprocity law (Chapter~10), and in modern language it says that a quaternion algebra over $\Q$ ramifies at an
even number of places.
\end{beauty}

\section{Hints}

Try each exercise for a while before reading its hint, and read the hint before the solution.

\begin{hint}{ex:fourrat}
Use Theorem~\ref{thm:localforms}(b) at each place, with $n=4$ and $n=3$.
\end{hint}
\begin{hint}{ex:fivesixths}
Count the $n\le x$ of the form $4^a(8b+7)$ for each $a$.
\end{hint}
\begin{hint}{ex:eureka}
$8\cdot\frac{k(k+1)}2+1=(2k+1)^2$. Write $8n+3$ as a sum of three squares and look modulo $8$.
\end{hint}
\begin{hint}{ex:legendre13}
Legendre's conditions with $(a,b,c)=(13,7,-1)$. Which squares are there modulo $7$ and modulo $13$?
\end{hint}
\begin{hint}{ex:prodQR}
Read the list in the proof of Theorem~\ref{thm:hilbertproduct} backwards.
\end{hint}
\begin{hint}{ex:x2y3z}
Only $p=2$ and $p=3$ can matter. At $3$, split according to $v_3(n)$; at $2$, use Theorem~\ref{thm:localforms}(a) for
$\langle1,2,3,-n\rangle$.
\end{hint}
\begin{hint}{ex:selmer}
In which $\Z_p$ is every unit a cube? For $p=3$, which units are cubes?
\end{hint}

\section{Solutions}

\begin{solution}{ex:fourrat}
\idea Representing $r$ by a form of rank $\ge4$ is automatic at every finite place, and rank $3$ fails at a place only if
$r$ is in one exceptional class.
\writeup Over $\R$, $\langle1,1,1,1\rangle$ and $\langle1,1,1\rangle$ represent every $r>0$. Over $\Q_p$, a form of rank $4$ represents
everything (Theorem~\ref{thm:localforms}(b)), so by Corollary~\ref{cor:hmrep} every $r>0$ is a sum of four rational squares (this
is weaker than Lagrange, but proved with no work). For $\langle1,1,1\rangle$: $d=1$, $\varepsilon=1$; it represents $r$ unless $r=-d=-1$
modulo squares and $(-1,-1)_p\ne\varepsilon=1$. At odd $p$, $(-1,-1)_p=1$, so there is no condition; at $p=2$, $(-1,-1)_2=-1$, so
$\langle1,1,1\rangle$ represents $r$ iff $r\ne-1$ modulo squares, i.e.\ iff $-r$ is not a square in $\Q_2$. Since $-7\equiv1\pmod8$,
$7$ is not a sum of three rational squares, nor is $\frac74$ or $28$ (same square class); but $14$ is ($-14$ has odd valuation).
\end{solution}

\begin{solution}{ex:fivesixths}
\idea The exceptions form disjoint progressions of density $\frac18,\frac1{32},\frac1{128},\dots$
\writeup For fixed $a$, the $n\le x$ with $n=4^a(8b+7)$ number $\frac x{8\cdot4^a}+O(1)$, and there are none for $4^a>x$. Summing over the
$O(\log x)$ values of $a$, the exceptions number $\frac x8\sum_a4^{-a}+O(\log x)=\frac x8\cdot\frac43+O(\log x)=\frac x6+O(\log x)$. So
the proportion of sums of three squares tends to $1-\frac16=\frac56$.
\end{solution}

\begin{solution}{ex:eureka}
\idea Triangular numbers are odd squares in disguise.
\writeup $8n+3\equiv3\pmod8$ is not of the form $4^a(8b+7)$, so $8n+3=x^2+y^2+z^2$. Squares are $0,1,4\pmod8$, and the only way to
get $3$ is $1+1+1$: $x,y,z$ are odd, $x=2k+1$ etc. Then $8n+3=\sum(8T_k+1)$ with $T_k=\frac{k(k+1)}2$, so $n=T_k+T_l+T_m$.
\end{solution}

\begin{solution}{ex:legendre13}
\idea Legendre's conditions are three square tests; the product formula guarantees that failures come in pairs.
\writeup With $a=13$, $b=7$, $c=-1$ (squarefree, coprime, mixed signs): we need $-bc=7$ to be a square modulo $13$, $-ca=13\equiv6$ a
square modulo $7$, and $-ab$ a square modulo $1$. The squares modulo $13$ are $1,4,9,3,12,10$, so $7$ is not one; the squares modulo
$7$ are $1,2,4$, so $6$ is not one. So there is no rational solution. In Hilbert symbols, $(-ac,-bc)_v=(13,7)_v$, which is
$\bigl(\frac7{13}\bigr)=-1$ at $13$, $\bigl(\frac{13}7\bigr)=-1$ at $7$, and $1$ at $2$ (since $\epsilon(13)=0$) and elsewhere:
two failures, as the product formula demands, and the equality $\bigl(\frac7{13}\bigr)=\bigl(\frac{13}7\bigr)$ is quadratic
reciprocity for $7$ and $13$.
\end{solution}

\begin{solution}{ex:prodQR}
\idea Each case of the product formula involves only two or three places, and the local formulas turn it into one of the three
laws.
\writeup For distinct odd primes $\ell,\ell'$: $1=\prod_v(\ell,\ell')_v=\bigl(\frac{\ell'}\ell\bigr)\bigl(\frac\ell{\ell'}\bigr)
(-1)^{\epsilon(\ell)\epsilon(\ell')}$, which is reciprocity. With $(-1,\ell)$: $1=\bigl(\frac{-1}\ell\bigr)(-1)^{\epsilon(\ell)}$, the
first supplement. With $(2,\ell)$: $1=\bigl(\frac2\ell\bigr)(-1)^{\omega(\ell)}$, the second supplement. (The local formulas were
proved by local arguments only, so there is no circularity.)
\end{solution}

\begin{solution}{ex:x2y3z}
\idea Apply Corollary~\ref{cor:hmrep}: only the places where the form has non-unit coefficients, $2$ and $3$, can obstruct; at $3$
there is always a way out, and at $2$ the obstruction is a single class.
\writeup Let $F=\langle1,2,3,-n\rangle$. Over $\R$ it is indefinite. At $p\ge5$, $\langle1,2,3\rangle$ is a unit form of rank $3$, isotropic.
At $p=3$: if $v_3(n)$ is even, $-n\sim-u$ with $u$ a unit and $\langle1,2,-u\rangle$ is a unit form of rank $3$, isotropic; if $v_3(n)$
is odd, $-n\sim-3u$, and either $u\equiv1\pmod3$ and $\langle3,-3u\rangle$ is isotropic, or $u\equiv2$ and $\langle1,2\rangle\cong
\langle1,-1\rangle$ ($-2\equiv1$ is a square in $\Q_3$) is isotropic. At $p=2$: by Theorem~\ref{thm:localforms}(a), $F$ is anisotropic
iff $d(F)=-6n$ is a square and $\varepsilon(F)\ne(-1,-1)_2=-1$. Writing $n=4^am$ with $4\nmid m$, $-6n$ is a square iff $m=2m'$ with
$-3m'\equiv1\pmod8$, i.e.\ $m'\equiv5\pmod8$, i.e.\ $n=4^a(16b+10)$. In that case $-n\sim6$ modulo squares and
\[
\varepsilon(F)=(2,3)(2,6)(3,6)=(2,2)(2,3)(3,3)=(2,-1)(2,3)(3,-1)=1\cdot(-1)\cdot(-1)=1,
\]
using $\omega(-1)=0$, $\omega(3)=1$, $\epsilon(3)=\epsilon(-1)=1$. So $F$ is anisotropic over $\Q_2$ exactly for $n=4^a(16b+10)$.
\end{solution}

\begin{solution}{ex:selmer}
\idea A cube root of a unit can be found by Hensel's lemma as soon as the derivative $3x^2$ is a unit, that is for $p\ne3$, and
cubing is a bijection on $\mathbb F_p^\times$ when $3\nmid p-1$. At $3$ one needs a congruence modulo $27$.
\writeup Over $\R$, take $y=z=1$ and $x=-\sqrt[3]3$. For $p=2,5$ and every $p\equiv2\pmod3$: cubing is a bijection of
$\mathbb F_p^\times$ (as $\gcd(3,p-1)=1$), and Hensel's lemma for $X^3-c$ at a root $a$ modulo $p$ ($|3a^2|_p=1$) shows that every
unit $c$ of $\Z_p$ is a cube. Take $z=0$, $y=1$ and $x^3=-\frac43$, a unit at $p\ne2,3$; for $p=2$ take instead $y=0$, $z=1$,
$x^3=-\frac53$, a $2$-adic unit. For $p=3$: take $x=0$, $z=1$, $y^3=-\frac54$. Modulo $27$, $4^{-1}\equiv7$, so $-\frac54\equiv-35
\equiv19$, and $7^3=343\equiv19\pmod{27}$. Hensel for $X^3-c$ at $a=7$: $|f(7)|_3\le3^{-3}<|f'(7)|_3^2=|147|_3^2=3^{-2}$; so $c$ is a
cube in $\Z_3$. For $p\equiv1\pmod3$, $p\ge7$, cubing is $3$-to-$1$ and the argument breaks down. One shows instead that the curve
$3x^3+4y^3+5z^3=0$ has a non-singular point modulo $p$, by counting points (Chapter~13: Hasse's bound gives at least $p+1-2\sqrt p>0$
points), and lifts it with Hensel's lemma. That the cubic has no rational point is Selmer's theorem (1951), much deeper.
\end{solution}

\section*{Sources and further reading}

\begin{itemize}
\item J.-P.~Serre, \emph{A Course in Arithmetic}, Springer GTM~7, 1973, Chapters~III and~IV. The route followed here, including the
existence theorem and the Davenport--Cassels lemma.
\item H.~Hasse, ``\"Uber die Darstellbarkeit von Zahlen durch quadratische Formen im K\"orper der rationalen Zahlen'', \emph{J.\
reine angew.\ Math.}\ 152 (1923), 129--148; and ``\"Uber die \"Aquivalenz quadratischer Formen im K\"orper der rationalen Zahlen'',
same journal 152 (1923), 205--224.
\item H.~Minkowski, ``Untersuchungen \"uber quadratische Formen'', \emph{Acta Math.}\ 7 (1885), and ``\"Uber die Bedingungen, unter
welchen zwei quadratische Formen mit rationalen Koeffizienten ineinander rational transformiert werden k\"onnen'', \emph{J.\ reine
angew.\ Math.}\ 106 (1890), 5--26.
\item J.~W.~S.~Cassels, \emph{Rational Quadratic Forms}, Academic Press, 1978 (Dover, 2008). Everything about forms over $\Q$ and
$\Z$, including Aubry's lemma.
\item E.~S.~Selmer, ``The Diophantine equation $ax^3+by^3+cz^3=0$'', \emph{Acta Math.}\ 85 (1951), 203--362.
\item K.~Ono and K.~Soundararajan, ``Ramanujan's ternary quadratic form'', \emph{Invent.\ Math.}\ 130 (1997), 415--454. What is
known (and conditional on GRH) about $x^2+y^2+10z^2$, the form Ramanujan could not finish.
\item \emph{History.} Gauss's diary entry of 10 July 1796: J.~J.~Gray, ``A commentary on Gauss's mathematical diary,
1796--1814'', \emph{Expo.\ Math.}\ 2 (1984), 97--130. Legendre's 1798 argument and its gap: A.~Weil, \emph{Number Theory: An
Approach through History}, Birkh\"auser, 1984, Chapter~IV.
\end{itemize}
